\subsection{De la generación de texto a la intervención en un repositorio}

Un agente de programación combina un modelo de lenguaje con herramientas para
observar y modificar un entorno de desarrollo. En este trabajo se utiliza el
término para describir un proceso que consulta archivos, propone modificaciones,
ejecuta comandos y utiliza sus resultados para continuar una tarea. Esta
definición operativa permite distinguir la propuesta de código de su comprobación:
una explicación coherente de una solución no demuestra que el programa resultante
satisfaga sus requisitos.

La documentación de Codex CLI describe una herramienta que inspecciona el
repositorio local, edita archivos y ejecuta las herramientas instaladas. También
contempla su utilización interactiva y mediante flujos repetibles con
\texttt{codex exec}, así como instrucciones reutilizables mediante
\textit{skills}~\cite{codex_cli_agents_2026}. Estas capacidades describen el
producto y no implican que todas ellas se hayan empleado en cada tarea de SCDE.
Conviene separar, además, el modelo que genera decisiones de la aplicación que
le proporciona acceso al sistema: el comportamiento observado depende de ambos
y de la configuración de la sesión.

Claude Code ofrece igualmente operaciones sobre varios archivos, ejecución de
pruebas y uso de Git. Su documentación presenta \texttt{CLAUDE.md} como contexto
persistente del proyecto y las \textit{skills} como una forma de agrupar
procedimientos reutilizables~\cite{codex_claude_overview_2026}. La coincidencia en
estas funciones permite analizar ambas herramientas como agentes que intervienen
en un ciclo de desarrollo. No constituye una comparación experimental entre
proveedores ni permite atribuirles una fiabilidad equivalente.

\subsection{Contexto explícito y coordinación del trabajo}

Las instrucciones de proyecto expresan restricciones que no se deducen de una
tarea aislada: estructura de directorios, contratos externos, convenciones de
interfaz y condiciones de entrega. La documentación oficial de Codex establece
que \texttt{AGENTS.md} proporciona instrucciones que se incorporan al inicio de
la ejecución y admite una organización por ámbitos del
proyecto~\cite{codex_agents_instructions_2026}. De ello se deriva una distinción
práctica: proporcionar una norma en el contexto orienta la generación, mientras
que impedir una escritura o rechazar una comprobación depende de mecanismos
externos al texto de esa norma.

En SCDE, \texttt{AGENTS.md}, \texttt{CLAUDE.md} y los documentos de coordinación
definen una distribución explícita de responsabilidades. Claude Code se asigna
a los flujos docentes y a la auditoría final; Codex, a administración, cuenta y
notificaciones. Cada línea utiliza una rama y un directorio de trabajo de Git
propios. \texttt{OWNERSHIP.json} enumera las rutas modificables y mantiene una
base compartida congelada; \texttt{REQUESTS.md} canaliza las solicitudes que
afectan a esa base~\cite{codex_scde_agent_evidence}. La separación limita el
solapamiento previsto, pero no demuestra que una integración posterior carezca
de conflictos funcionales.

El contrato de la API y la guía de interfaz cumplen funciones complementarias.
El primero acota los mensajes que puede enviar el frontend; la segunda fija
criterios de composición y presentación. Los registros de trabajo documentan
decisiones y restricciones encontradas, de modo que una tarea puede retomarse
sin tratar una suposición anterior como un hecho confirmado. En este proyecto,
esta trazabilidad es especialmente relevante porque el mismo proceso produce
código, pruebas, scripts de publicación y borradores de documentación.

\subsection{Evaluación y alcance de la evidencia}

SWE-bench plantea la resolución de incidencias reales de repositorios como tarea
de evaluación: el sistema recibe un repositorio y una descripción del problema,
y debe modificar el código para resolverlo. El artículo destaca que estas tareas
pueden exigir cambios coordinados entre funciones y archivos, además de
interacción con el entorno de ejecución~\cite{codex_swebench_2024}. Este marco
resulta pertinente para distinguir completar una función aislada de mantener un
sistema existente. Sus resultados no se trasladan automáticamente a un frontend
React ni a las versiones de herramientas utilizadas en SCDE.

La verificación local de SCDE combina análisis estático, pruebas automatizadas,
construcción y auditorías de estilo y propiedad. La inspección del auditor de
propiedad muestra un límite concreto: compara los commits de la rama con la
base y no incorpora por sí mismo los cambios todavía sin confirmar. Por ello,
su resultado debe interpretarse junto con la revisión de los archivos que se
pretende confirmar~\cite{codex_scde_agent_evidence}. Del mismo modo, una prueba
que utiliza respuestas simuladas aporta evidencia sobre ese escenario y no
certifica el comportamiento del backend Django.

Los fallos de permisos o de ejecución de procesos deben registrarse como
restricciones de la sesión. No equivalen a un defecto funcional demostrado,
pero tampoco permiten declarar superada la comprobación que quedó interrumpida.
La bitácora distingue estas situaciones de los errores reproducidos por las
pruebas. Sin un protocolo comparativo, una línea base y medidas homogéneas de
tiempo y revisión, este trabajo no atribuye a los agentes un porcentaje de
productividad ni una reducción cuantificada de defectos.

\subsection{Responsabilidad y transparencia de la asistencia}

Claude Code y Codex se han utilizado, bajo la dirección y las reglas establecidas
por el autor, para elaborar código y borradores de la memoria. Parte del trabajo
se ha ejecutado mediante sesiones autónomas: la delegación de tareas no debe
presentarse como revisión humana continua de cada cambio. Las comprobaciones
automáticas, las bitácoras y la revisión por otro agente aportan evidencia
complementaria, pero la revisión y aceptación final del autor siguen siendo
necesarias~\cite{codex_scde_agent_evidence}.

Esta sección sitúa las herramientas y los límites de su evaluación. El detalle
de las tareas automatizadas, los procedimientos de comprobación y la revisión
del material generado debe quedar desarrollado también en la metodología y en
la declaración de asistencia de la memoria. La revisión entre agentes no debe
describirse como validación humana independiente ni sustituir la comprobación
de las fuentes bibliográficas.
